\documentclass[12pt]{article}
\usepackage[utf8]{inputenc}
\usepackage[spanish]{babel}
\usepackage[a4paper, margin=2cm]{geometry}
\usepackage{tikz}
\usetikzlibrary{babel}
\usepackage[most]{tcolorbox}
\usepackage{amsmath}

\renewcommand{\familydefault}{\sfdefault}

% Definición de colores
\definecolor{medblue}{RGB}{43, 108, 176}
\definecolor{medteal}{RGB}{49, 151, 149}
\definecolor{medlight}{RGB}{235, 248, 255}
\definecolor{meddark}{RGB}{26, 32, 44}
\definecolor{skin}{RGB}{253, 218, 184}        % Color piel
\definecolor{skindark}{RGB}{203, 143, 98}     % Borde color piel
\definecolor{salivablue}{RGB}{147, 197, 253}  % Azul claro para gotas/aire

\begin{document}
\pagestyle{empty}

\begin{tcolorbox}[
    enhanced, 
    colback=medlight!30, 
    colframe=medblue, 
    title={\Large \textbf{Energía Cinética: Reflejo del Estornudo}}, 
    halign title=center, 
    fonttitle=\bfseries, 
    arc=4mm, 
    boxrule=1.5pt,
    shadow={2mm}{-2mm}{0mm}{black!30!white}
]

    \vspace{0.2cm}
    \begin{tcolorbox}[colback=white, colframe=medteal, arc=2mm, boxrule=1.2pt, title=\textbf{Caso Clínico / Problema}]
        Durante un estornudo, una persona expulsa aproximadamente \textbf{0.5 g} (\textbf{0.0005 kg}) de aire y pequeñas gotas de saliva a una velocidad asombrosa de \textbf{40 m/s}. Calcule la energía cinética de esta masa expulsada.
    \end{tcolorbox}

    \vspace{0.4cm}
    
    \begin{center}
    \begin{tikzpicture}[scale=0.95]
        
        % --- Silueta de la cabeza (Perfil izquierdo) ---
        \filldraw[skin, draw=skindark, thick]
            (-1, -2) .. controls (-1.5, -1) and (-1.5, 1) .. (0, 2)    % Nuca a coronilla
            .. controls (1, 2.5) and (2, 1.5) .. (2, 0.5)              % Frente
            .. controls (2.2, 0.2) and (2.5, -0.1) .. (2.3, -0.3)      % Nariz
            .. controls (2.1, -0.4) and (2.0, -0.5) .. (2.1, -0.5)     % Labio superior
            -- (1.0, -0.8)                                             % Interior de la boca (abierta)
            -- (2.1, -1.1)                                             % Labio inferior
            .. controls (1.8, -1.3) and (1.5, -1.8) .. (1.5, -2)       % Cuello
            -- cycle;
            
        % Ojo estilizado
        \filldraw[meddark] (1.3, 0.5) ellipse (0.15cm and 0.08cm);
        
        % --- Nube de estornudo (Cono de aire) ---
        \fill[salivablue, opacity=0.3] (1.5, -0.8) -- (7, 1.5) .. controls (8, -0.8) and (7, -3) .. (7, -3) -- cycle;
        
        % Gotas de saliva y partículas
        \foreach \x/\y in {3/0.2, 3.5/-0.5, 4/0.5, 4.5/-1, 5/0.8, 5.5/-0.2, 6/0.4, 6.5/-1.5, 6.2/-0.8, 4.8/-1.8} {
            \filldraw[salivablue!80!blue, draw=meddark, thin] (\x, \y) circle (0.1cm);
            \filldraw[salivablue!80!blue, draw=meddark, thin] (\x-0.3, \y+0.3) circle (0.05cm);
        }
        
        % --- Vectores e Información ---
        \draw[->, ultra thick, red!80!black] (2.5, -0.8) -- (6.0, -0.8) node[midway, above, font=\large\bfseries, fill=white, inner sep=2pt, rounded corners=2pt, opacity=0.9, text opacity=1] {Velocidad $v = 40 \text{ m/s}$};
        
        \node[font=\bfseries, text=meddark, fill=white, inner sep=3pt, rounded corners=3pt, draw=meddark!50, thick] at (4.5, -2.5) {Masa expulsada $m = 0.0005 \text{ kg}$};
        
    \end{tikzpicture}
    \end{center}

    \vspace{0.4cm}
    \begin{tcolorbox}[colback=medteal!10, colframe=medteal, arc=2mm, boxrule=1.2pt, title=\textbf{Desarrollo Matemático y Solución}]
        La energía cinética ($E_c$) es la energía asociada al movimiento de una masa. Se calcula mediante la fórmula $E_c = \frac{1}{2} m v^2$. Es crucial usar la masa en kilogramos para que el resultado esté en Joules (J).
        \begin{align*}
            E_c &= \frac{1}{2} \cdot (0.0005 \text{ kg}) \cdot (40 \text{ m/s})^2 \\[1ex]
            E_c &= 0.00025 \text{ kg} \cdot 1600 \text{ m}^2/\text{s}^2 \\[1ex]
            E_c &= \mathbf{0.4 \text{ Joules}}
        \end{align*}
        \textbf{Resultado:} La energía cinética liberada por las partículas durante este estornudo es de \textbf{0.4 Joules}.
    \end{tcolorbox}
    
\end{tcolorbox}

\end{document}